\documentclass[10pt]{article}
\usepackage[margin=0.8in]{geometry}
\usepackage[T1]{fontenc}
\usepackage{lmodern}
\usepackage{amsmath,amssymb,booktabs,array}
\usepackage{graphicx}
\usepackage{xcolor}
\usepackage{enumitem}
\usepackage{microtype}
\usepackage{hyperref}
\usepackage{titlesec}
\definecolor{accent}{HTML}{0F766E}
\hypersetup{colorlinks=true,linkcolor=accent,urlcolor=accent}
\setlength{\parindent}{0pt}
\setlength{\parskip}{5pt}
\titlespacing*{\section}{0pt}{14pt}{5pt}
\titlespacing*{\subsection}{0pt}{9pt}{3pt}
\newcommand{\code}[1]{\texttt{#1}}
\newcommand{\E}{\mathbb{E}}
\title{\vspace{-1.2em}\textbf{Fractal Potential-Field RL}\\
\large Physics Revision and Experimental Audit}
\author{Technical implementation note}
\date{September 2026}
\begin{document}
\maketitle
\vspace{-1.2em}
\begin{abstract}
This report audits the replacement of resistance-dependent drag by a conservative
fractal potential field and linear damping. The earlier weak-field result is not
used as evidence: fine terrain and low damping let the agent coast through it.
The audited valley-field run uses broad low-frequency basins, calibrated field
acceleration, stronger damping, and a small control resource cost. It documents
the physical law, PPO objective, 3090 result, potential-overlaid trajectory, and
the limitation that speed reward alone does not require exclusive valley following.
\end{abstract}

\section{Problem formulation}
The environment is a continuous 2-D agent moving on a periodic $512\times512$
pink-noise potential map. Positions are unbounded in physical coordinates, while
terrain queries use periodic bilinear sampling. Intensity $I$ is generated by
Fourier-filtered white Gaussian noise; its finite-sample mean is subtracted and
the physical potential is $\phi=(I-\bar I)^2$. The clean
audit uses spectral exponent $\beta=3$, intentionally making broad basins and
connected valleys rather than pixel-scale corrugations. The simulator uses
$\Delta t=0.02$ with two substeps.

At time $t$, the agent state comprises position $\mathbf{x}_t$, velocity
$\mathbf{v}_t$, heading, mass, and the prior action. The policy chooses
$\mathbf{a}_t=(a_{\parallel,t},a_{\perp,t})\in[-1,1]^2$. These components command
signed acceleration along the current heading and to its right, with configured
maximum magnitudes 2.0 and 1.0. Near zero speed, the stored heading defines the
control frame. The field contribution is conservative and the resistance is
linear in velocity. The implemented acceleration is
\begin{equation}
 \mathbf a=\mathbf a_{\rm ctl}-s_\phi\frac{\nabla\phi}{g_{95}}-\eta\mathbf v,
\end{equation}
where $g_{95}$ is the stored 95th percentile of $\|\nabla\phi\|$. Thus
$s_\phi=2.0$ is a terrain acceleration comparable directly to maximum forward
control acceleration 2.0, and $\mathbf a_{\rm resist}=-\eta\mathbf v$ is the
requested linear resistance acceleration.

\begin{center}
\begin{tabular}{lll}
\toprule
Quantity & Audited value & Interpretation\\
\midrule
Spectral exponent $\beta$ & 3.0 & broad valley geometry\\
Potential scale $s_\phi$ & 2.0 & 95th-percentile field acceleration\\
Forward control limit & 2.0 & bounded control acceleration\\
Linear damping $\eta$ & 0.6 & $v_\infty=2/0.6=3.33$ on flat terrain\\
Control penalty & 0.05 & explicit acceleration resource cost\\
\bottomrule
\end{tabular}
\end{center}

The per-step reward is
\begin{equation}
  r_t = w_s \, \|\mathbf{v}_t\|\Delta t
  - \lambda_a\|\mathbf{a}_t\|^2
  - \lambda_{\Delta a}\|\mathbf{a}_t-\mathbf{a}_{t-1}\|^2,
\end{equation}
where $w_s=1$, the audited control penalty is 0.05, and action-change penalty is
zero. Thus the task
maximizes travelled distance, equivalently sustained mean speed over the fixed
2000-step episode, rather than displacement toward a goal. Time-limit endings
are truncations, not physical terminal states.

\section{Observation and policy interface}
\subsection{Inputs}
The observation intentionally omits absolute position and terrain seed. It has:
\begin{itemize}[leftmargin=1.4em,itemsep=2pt]
  \item a one-channel $32\times32$ potential patch, centered on the body and
  rotated into the direction of travel (the image top is forward); and
  \item a six-element state vector: local forward velocity, local right velocity,
  speed, mass, previous throttle, and previous steering.
\end{itemize}
This makes the task locally observable: the policy must infer useful motion
decisions from terrain immediately around the vehicle and its current dynamical state.

\section{Power-limited mode and randomized tiled starts}
The default controller remains acceleration-limited. An optional configuration,
\code{physics.power\_limit}, additionally caps longitudinal mechanical power:
\begin{equation}
 |a_{\parallel}|\leq\min\left(a_{\parallel,\max},
 \frac{P_{\max}}{m\max(\|\mathbf v\|,\epsilon)}\right).
\end{equation}
The small $\epsilon$ retains the ordinary launch acceleration cap at rest. The
cap applies to the magnitude of longitudinal thrust or braking; lateral control
is unchanged. The supplied \code{config/power\_limited.yaml} sets $P_{\max}=2$.
At every reset, \code{environment.random\_start\_position=true} samples a fresh,
uniform location in the fundamental periodic tile, avoiding a fixed starting
location while preserving the same tiled landscape.

\begin{figure}[t]
\centering
\includegraphics[width=\linewidth]{dumps/plots/power_limited_mode.png}
\caption{New diagnostics. Left: potential made by centering the filtered Gaussian
field and squaring it. Right: the optional power cap retains full acceleration at
low speed and throttles the magnitude of longitudinal acceleration as speed rises.}
\label{fig:power-limit}
\end{figure}

\subsection{Shared convolutional encoder}
The actor and critic use one shared feature extractor. For the default field of
view, the CNN maps $1\times32\times32$ terrain pixels to a 1152-element vector.
After concatenating the six scalar features, a two-layer MLP creates the 128-D
latent representation $\mathbf{z}_t$. This is the representation dumped by the
new analysis tool.

\begin{center}
\renewcommand{\arraystretch}{1.22}
\begin{tabular}{>{\raggedright\arraybackslash}p{0.26\linewidth}p{0.48\linewidth}p{0.16\linewidth}}
\toprule
\textbf{Stage} & \textbf{Operation} & \textbf{Output} \\
\midrule
Terrain input & potential patch & $1\times32\times32$ \\
Conv block 1 & Conv2d($1\rightarrow16$, kernel 5, stride 2), ReLU & $16\times14\times14$ \\
Conv block 2 & Conv2d($16\rightarrow32$, kernel 3, stride 2), ReLU & $32\times6\times6$ \\
Flatten & spatial feature flattening & 1152 \\
Fusion & concatenate six state values & 1158 \\
Encoder MLP 1 & Linear($1158\rightarrow128$), Tanh & 128 \\
Encoder MLP 2 & Linear($128\rightarrow128$), Tanh & 128 \\
Actor & Linear($128\rightarrow2$), learned log standard deviation & 2 means \\
Critic & Linear($128\rightarrow1$) & 1 value \\
\bottomrule
\end{tabular}
\end{center}

Concretely,
\begin{equation}
 \mathbf{z}_t =
 \operatorname{Tanh}\left(W_2
 \operatorname{Tanh}\left(W_1[\operatorname{CNN}(I_t);\mathbf{s}_t]+b_1\right)+b_2\right).
\end{equation}
The actor predicts $\boldsymbol{\mu}_t=W_\pi\mathbf{z}_t+b_\pi$ and the critic
predicts $V_\phi(o_t)=W_V\mathbf{z}_t+b_V$. There is one trainable log standard
deviation per action dimension, initialized at $-0.5$ and clamped to $[-5,2]$
before exponentiation. The sampled raw action is
$\mathbf{u}_t\sim\mathcal{N}(\boldsymbol{\mu}_t,\operatorname{diag}(\boldsymbol{\sigma}^2))$,
and the executed action is $\mathbf{a}_t=\tanh(\mathbf{u}_t)$. The implementation
uses the exact tanh change-of-variables correction in the action log probability.

\section{PPO optimization}
\subsection{Rollouts and bootstrapped returns}
The default trainer runs 16 environments synchronously. Each rollout has 128
simulator steps, yielding $16\times128=2048$ transitions per update. Let
$V_t=V_\phi(o_t)$ and $V_{t+1}$ denote the critic evaluated on the actual next
observation. The temporal-difference residual is
\begin{equation}
 \delta_t=r_t+\gamma(1-d_t)V_{t+1}-V_t,
\end{equation}
where $d_t$ is one for true termination. Generalized advantage estimation (GAE)
uses
\begin{equation}
 \hat A_t=\delta_t+\gamma\lambda(1-e_t)\hat A_{t+1},
 \qquad \hat R_t=\hat A_t+V_t,
\end{equation}
where $e_t$ is one for any episode ending. This distinction is deliberate:
time-limit truncation bootstraps from $V_{t+1}$, but the advantage recursion does
not leak across the subsequent reset. Advantages are normalized within each batch.

\subsection{Clipped policy and value losses}
For data gathered by the old policy, define
$r_t(\theta)=\exp(\log\pi_\theta(\mathbf{a}_t|o_t)-\log\pi_{\rm old}(\mathbf{a}_t|o_t))$.
The policy loss is the standard clipped surrogate,
\begin{equation}
 L_\pi(\theta)=-\E_t\left[\min\left(r_t(\theta)\hat A_t,
 \operatorname{clip}(r_t(\theta),1-\epsilon,1+\epsilon)\hat A_t\right)\right],
\end{equation}
and the critic loss is
\begin{equation}
 L_V(\phi)=\frac12\E_t[(V_\phi(o_t)-\hat R_t)^2].
\end{equation}
The implementation minimizes
\begin{equation}
 L=L_\pi+c_VL_V-c_H\E_t[H(\pi_\theta(\cdot|o_t))].
\end{equation}
For this run, $\gamma=0.99$, $\lambda=0.95$, $\epsilon=0.2$,
$c_V=0.5$, $c_H=0$, learning rate $3\cdot10^{-4}$, and maximum gradient norm
0.5. Adam performs four epochs over each rollout, split into minibatches of 256.
The entropy term is implemented for completeness but inactive in the default
configuration because $c_H=0$.

\section{JEPA auxiliary objective}
The reported 3090 experiments optimize PPO only. The implementation now exposes
a configurable Joint-Embedding Predictive Architecture (JEPA) auxiliary objective so
the actor encoder is trained to predict future compact representations rather
than reconstructing terrain pixels. Let $f_\theta(o_t)=\mathbf z_t$ be the online
encoder, $\bar f(o_{t+k})=\bar{\mathbf z}_{t+k}$ be a stop-gradient target
encoder updated by exponential moving average, and $q_\psi$ be a predictor that
receives the current embedding and intervening actions. For a prediction horizon
$k\in\{1,\ldots,K\}$,
\begin{equation}
 \mathcal L_{\rm JEPA}=\frac{1}{K}\sum_{k=1}^{K}
 \left\|q_\psi\left(\mathbf z_t,\mathbf a_{t:t+k-1},k\right)
 -\operatorname{sg}[\bar{\mathbf z}_{t+k}]\right\|_2^2,
 \qquad
 \bar\theta\leftarrow\tau\bar\theta+(1-\tau)\theta.
\end{equation}
Here $\operatorname{sg}$ stops gradients through the target branch. The implemented
joint optimization is $L_{\rm total}=L_{\rm PPO}+\lambda_{\rm JEPA}\mathcal
L_{\rm JEPA}$, with predictor and online encoder updated by gradients and the
target encoder updated only by EMA. This makes representation changes measurable
with the pairwise UMAP/CKA pipeline while preserving the PPO policy objective.
It is available through \code{training.jepa\_coef} and was \emph{not} used to produce
the current reported curves, UMAPs, or CKA matrix.

\section{New-model test: protocol and results}
The training bank contains 64 terrain seeds beginning at the configured seed.
Evaluation uses the disjoint fixed seeds 100001--100004. The current test uses
the standardized centered-Gaussian potential, $P_{\max}=6$ power cap, and
unshared value-clipped PPO for 10,000 transitions on an NVIDIA RTX 3090. Every
figure below is regenerated from this checkpoint and held-out inference.

\begin{figure}[t]
\centering
\includegraphics[width=\linewidth]{dumps/power_limited_6_10k/plots/learning_curve.png}
\caption{Corrected RTX 3090 run: rollout speed (gray) and deterministic
held-out mean speed (teal). The band is standard deviation over four held-out seeds.}
\label{fig:learning}
\end{figure}

Figure~\ref{fig:learning} is a smoke experiment, not a convergence claim. The
corrected RTX 3090 run reaches held-out mean speed 0.818 at 10,000 transitions,
after climbing from 0.055 at the first checkpoint. It validates the corrected
terrain scale, power-limited controller, and PPO pipeline end-to-end.

\begin{center}
\begin{tabular}{rrrr}
\toprule
\textbf{Transitions} & \textbf{Held-out mean speed} & \textbf{Std.} & \textbf{Interpretation}\\
\midrule
2,048 & 0.055 & 0.035 & corrected 3090 initial checkpoint\\
6,144 & 0.335 & 0.034 & corrected 3090 intermediate checkpoint\\
10,000 & 0.818 & 0.116 & corrected 3090 endpoint\\
\bottomrule
\end{tabular}
\end{center}

\subsection{3090 corrected power-limited rerun}
After changing the terrain definition to the centered-Gaussian squared potential,
standardizing the centered field before squaring it, and enabling the supplied
power-limited configuration, a fresh 10,000-transition run was executed on an
NVIDIA RTX 3090. It used $P_{\max}=6$, the same 16-environment
rollout protocol, and held-out terrain seeds 100001--100004. The held-out mean
speed was 0.818 at the 10,000-transition endpoint (intermediate values: 0.055,
0.134, 0.335, and 0.541 at 2,048, 4,096, 6,144, and 8,192 transitions).

This is an execution result, not a converged-performance claim. It is materially
better than the preliminary $P_{\max}=2$ run, which did not learn. Standardizing
the Gaussian potential restored a stable policy input range, while $P_{\max}=6$
begins throttling only near speed 3. A matched sweep of $P_{\max}$ values and a
longer run are required before choosing a final cap.

\section{Latent trajectory dumps and UMAP}
The analysis command \code{python -m driver.analyze\_latent CHECKPOINT --device cuda}
runs the deterministic policy on a held-out terrain. At every step it records
physical position, speed, potential, action, critic value, and the exact
128-D output of \code{ActorCritic.encode}. The raw trajectory is stored in
\code{dumps/latents/} as both a compressed NumPy archive and a CSV with
\code{latent\_000} through \code{latent\_127} columns. This is a dump of the
shared representation before either actor or critic head, not a visualization
of the two action outputs.

For visualization, each latent coordinate is standardized across the trajectory.
For cross-checkpoint comparison, each encoder is applied to the same 1000-step
held-out observation trajectory. Every row/column pair is jointly embedded with
UMAP using Euclidean distance, \code{n\_neighbors=30}, \code{min\_dist=0.1}, and
the training seed. The resulting 5x5 grid is shown in Figure~\ref{fig:umap}.
UMAP is a nonlinear local-neighborhood visualization: nearby points indicate
similar encoder activations, whereas absolute axes and long-range distances are
not calibrated physical distances.

The physical trajectory is additionally color-coded by the two UMAP coordinates:
normalized UMAP-1 is the red channel and normalized UMAP-2 is the blue channel.
This produces a direct visual correspondence between a location on the map and
the encoder representation without treating either UMAP axis as a physical
coordinate. Black denotes locations where both normalized coordinates are low.

\begin{figure}[t]
\centering
\includegraphics[width=\linewidth]{dumps/power_limited_6_pairwise_10k/plots/pairwise_encoder_umap_5x5_seed100001.png}
\caption{Corrected 3090 pairwise joint UMAPs across checkpoints 2,048, 4,096,
6,144, 8,192, and 10,000, using the same held-out inference trajectory. Blue is
the row encoder and orange is the column counterpart; diagonal cells compare an
encoder with itself.}
\label{fig:umap}
\end{figure}

\subsection{Velocity-frame acceleration diagnostics}
For every analyzed checkpoint, the same deterministic rollout also records the
signed components of the policy input acceleration and potential acceleration in
the instantaneous velocity frame. With $\hat{\mathbf f}=\mathbf v/\|\mathbf v\|$
at nonzero speed (and stored heading at rest) and
$\hat{\mathbf l}=(-\hat f_y,\hat f_x)$, the diagnostic reports
\begin{equation}
 a_{\parallel}=\mathbf a\mathbin{\cdot}\hat{\mathbf f},\qquad
 a_{\perp}=\mathbf a\mathbin{\cdot}\hat{\mathbf l}.
\end{equation}
Positive perpendicular acceleration is to the agent's right in new experiments.
Legacy checkpoints without an explicit convention retain their original
left-positive semantics. This separates
the controller's intended longitudinal/lateral input from the terrain's
longitudinal/lateral response; it is more informative than a magnitude-only
acceleration trace.

\begin{figure}[t]
\centering
\includegraphics[width=0.76\linewidth]{dumps/power_limited_6_pairwise_10k/plots/encoder_cka_matrix_seed100001.png}
\caption{Linear CKA matrix for the same corrected 3090 checkpoints and held-out
trajectory. Unlike joint UMAP, CKA is quantitative and remains invariant to
feature-coordinate rotations and permutations.}
\label{fig:episode-diagnostic}
\end{figure}

\section{Appended run: SAC-only 3090 comparison (September 20, 2026)}
This independent run uses the same standardized centered-Gaussian terrain,
$P_{\max}=6$ controller, 16-environment collection, held-out seeds, and
10,000-transition budget as the PPO smoke run. SAC uses a replay buffer,
independent actor encoder, twin action-value critics, target critics with
$\tau=0.005$, and learned entropy temperature. JEPA was disabled
($\lambda_{\rm JEPA}=0$), so this is a clean PPO-versus-SAC baseline rather
than a representation-learning ablation.

The held-out mean speed rose from 0.064 at 2,048 transitions to 1.655 at 10,000
transitions (standard deviation 0.056 across four held-out seeds). This is a
short-run comparison, not a convergence claim. It is appended as a separate
result; no PPO figures or metrics above were replaced.

\begin{figure}[t]
\centering
\includegraphics[width=\linewidth]{dumps/sac_10k/plots/learning_curve.png}
\caption{Appended SAC-only RTX 3090 run: rollout and deterministic held-out
mean speed. The lower panel records JEPA loss, which is zero because the
auxiliary objective is disabled for this baseline.}
\label{fig:sac-learning}
\end{figure}

\begin{figure}[t]
\centering
\includegraphics[width=\linewidth]{dumps/sac_10k/plots/latent_umap_sac_000010000_seed100001.png}
\caption{Appended SAC endpoint deterministic inference on held-out seed 100001.
Left: trajectory over the terrain-potential overlay. Right: the SAC actor
encoder trajectory in UMAP, colored by physical speed.}
\label{fig:sac-trajectory-umap}
\end{figure}

\begin{figure}[t]
\centering
\includegraphics[width=\linewidth]{dumps/sac_10k/plots/dynamics_history_sac_000010000_seed100001.png}
\caption{Appended SAC full dynamics diagnostic: physical and latent velocity,
controls, control/field/damping/net acceleration magnitudes, and signed
parallel/perpendicular acceleration over the same terrain-overlay inference.}
\label{fig:sac-dynamics}
\end{figure}

\begin{figure}[t]
\centering
\includegraphics[width=\linewidth]{dumps/sac_10k/plots/pairwise_encoder_umap_5x5_seed100001.png}
\caption{Appended SAC pairwise joint UMAP grid across checkpoints 2,048 through
10,000 on the same held-out probe trajectory.}
\label{fig:sac-pairwise}
\end{figure}

\begin{figure}[t]
\centering
\includegraphics[width=0.76\linewidth]{dumps/sac_10k/plots/encoder_cka_matrix_seed100001.png}
\caption{Appended SAC linear CKA matrix across the same checkpoints.}
\label{fig:sac-cka}
\end{figure}

\section{Appended run: PPO+JEPA 3090 validation (September 20, 2026)}
This is a new, isolated 10,000-transition RTX 3090 run, appended without
altering the earlier PPO-only or SAC-only evidence. It uses the same
standardized centered-Gaussian terrain, periodic tiling, random starts,
$P_{\max}=6$ longitudinal power cap, 16 parallel environments, and held-out
terrain seeds as the PPO baseline. The only learning change is the
action-conditioned one-step JEPA auxiliary loss with
\code{training.jepa\_coef=0.1} and EMA target update
\code{training.jepa\_target\_tau=0.01}.

Concretely, the implemented branch predicts the next target representation
from the current actor representation and the executed action:
\begin{equation}
 \mathcal L_{\rm JEPA}=\left\|q_\psi( f_\theta(o_t),a_t)
 -\operatorname{sg}[\bar f(o_{t+1})]\right\|_2^2,
 \qquad \bar\theta\leftarrow\tau\bar\theta+(1-\tau)\theta.
\end{equation}
The online actor encoder and predictor receive gradients; the target encoder
is an EMA copy. This is the implemented one-step instance of the more general
future-prediction objective introduced above. The stored transition uses the
actual pre-reset next observation, so a time-limit truncation does not corrupt
the JEPA target.

The held-out mean speed rose from $0.056\pm0.034$ at 2,048 transitions to
$0.676\pm0.105$ at 10,000. The auxiliary loss remained finite throughout and
decreased from 0.003337 to 0.000595 (an 82\% reduction), confirming that the
predictor/EMA path is active rather than merely configured.

\begin{figure}[t]
\centering
\includegraphics[width=\linewidth]{dumps/ppo_jepa_10k/plots/learning_curve.png}
\caption{Appended PPO+JEPA RTX 3090 learning curve. Top: rollout and
deterministic held-out speed. Bottom: finite JEPA prediction loss, decreasing
over the five saved checkpoints.}
\label{fig:jepa-learning}
\end{figure}

\begin{figure}[t]
\centering
\includegraphics[width=\linewidth]{dumps/ppo_jepa_10k/plots/latent_umap_ppo_000010000_seed100001.png}
\caption{Appended PPO+JEPA endpoint inference on held-out seed 100001. Left:
terrain-potential overlay and deterministic trajectory. Right: actor-encoder
UMAP colored by physical speed.}
\label{fig:jepa-trajectory-umap}
\end{figure}

\begin{figure}[t]
\centering
\includegraphics[width=\linewidth]{dumps/ppo_jepa_10k/plots/dynamics_history_ppo_000010000_seed100001.png}
\caption{Appended PPO+JEPA full dynamics diagnostic: physical and latent
velocity, control commands, control/field/damping/net accelerations, and
signed parallel/perpendicular accelerations for the terrain-overlay episode.}
\label{fig:jepa-dynamics}
\end{figure}

\begin{figure}[t]
\centering
\includegraphics[width=\linewidth]{dumps/ppo_jepa_10k/plots/pairwise_encoder_umap_5x5_seed100001.png}
\caption{Appended PPO+JEPA pairwise joint UMAP grid across checkpoints 2,048
through 10,000 on one common held-out inference trajectory.}
\label{fig:jepa-pairwise}
\end{figure}

\begin{figure}[t]
\centering
\includegraphics[width=0.76\linewidth]{dumps/ppo_jepa_10k/plots/encoder_cka_matrix_seed100001.png}
\caption{Appended PPO+JEPA linear CKA matrix for the same encoder snapshots.}
\label{fig:jepa-cka}
\end{figure}

\subsection{What differs across the 10k-transition 3090 runs}
\begin{center}
\begin{tabular}{>{\raggedright\arraybackslash}p{0.20\linewidth}>{\raggedright\arraybackslash}p{0.22\linewidth}>{\raggedright\arraybackslash}p{0.18\linewidth}>{\raggedright\arraybackslash}p{0.28\linewidth}}
\toprule
\textbf{Run} & \textbf{Learning setup} & \textbf{Held-out speed at 10k} & \textbf{Reading}\tabularnewline
\midrule
PPO-only & PPO, $\lambda_{\rm JEPA}=0$ & $0.818\pm0.116$ & Reference on-policy baseline under the same terrain and power cap.\tabularnewline
PPO+JEPA & PPO, one-step JEPA, $\lambda_{\rm JEPA}=0.1$ & $0.676\pm0.105$ & JEPA is operating and learns a stable predictive objective, but trails PPO-only by 0.142 (17.3\%) in this short single-run comparison.\tabularnewline
SAC-only & Replay SAC, twin Q critics, $\lambda_{\rm JEPA}=0$ & $1.655\pm0.056$ & Strongest endpoint here; this is an algorithm change, not evidence for a JEPA effect.\tabularnewline
SAC+JEPA & Replay SAC, one-step JEPA, $\lambda_{\rm JEPA}=0.1$ & $1.278\pm0.042$ & Active auxiliary branch, but lower than SAC-only in the initial short run; see the appended cadence erratum.\tabularnewline
\bottomrule
\end{tabular}
\end{center}

The immediate conclusion is deliberately narrow: the PPO+JEPA branch works
numerically and produces the complete inference/representation diagnostics,
but its current coefficient and 10k budget do not improve speed over PPO-only.
SAC reaches about $2.0\times$ the PPO-only endpoint and $2.45\times$ the
PPO+JEPA endpoint, but off-policy replay, twin critics, and entropy tuning are
confounded with that difference. A fair JEPA conclusion needs matched
multi-seed, longer-horizon PPO sweeps, beginning with lower coefficients such
as 0.025 and 0.05. SAC+JEPA was unimplemented at the time of the original
comparison and is now reported in the appended section below.

\section{Appended run: SAC+JEPA 3090 validation (September 20, 2026)}
This section appends the now-implemented SAC+JEPA result; it does not revise
the preceding report, whose statement about the then-unimplemented branch was
correct when written. The SAC actor now owns an online terrain/state encoder,
an EMA target copy of that encoder, and a 130-to-128 predictor. Replay supplies
the executed tanh-squashed action and the real pre-reset next observation, so
the actor receives
\begin{equation}
 \mathcal L_{\rm actor}^{\rm SAC+JEPA}=
 \mathbb E[\alpha\log\pi(a|o)-\min(Q_1,Q_2)]
 +0.1\left\|q_\psi([f_\theta(o_t),a_t])-
 \operatorname{sg}[\bar f(o_{t+1})]\right\|_2^2 .
\end{equation}
The JEPA target update rate is 0.01 and is distinct from SAC's target-critic
rate $\tau=0.005$. Target-encoder/predictor state is stored in the SAC actor
checkpoint, and standard inference loading restores it.

This isolated RTX 3090 run uses the same centered-Gaussian terrain, random
starts, $P_{\max}=6$ cap, 16-environment collection, replay configuration,
and held-out terrain seeds as the SAC-only baseline. At 10,000 transitions,
SAC+JEPA reaches held-out mean speed $1.278\pm0.042$; its nonzero auxiliary
loss is $9.08\times10^{-6}$, $7.22\times10^{-6}$, $5.80\times10^{-6}$,
$2.25\times10^{-5}$, and $3.52\times10^{-5}$ across the five checkpoints.
Thus the branch is active, but the very low scale and late rise make it a
representation-shaping signal rather than evidence of a useful predictive
improvement at this coefficient and horizon.

\begin{figure}[t]
\centering
\includegraphics[width=\linewidth]{dumps/sac_jepa_10k_v2/plots/learning_curve.png}
\caption{Appended SAC+JEPA RTX 3090 learning curve. The lower panel reports
the actual replay-based JEPA loss, not a hardcoded zero.}
\label{fig:sac-jepa-learning}
\end{figure}

\begin{figure}[t]
\centering
\includegraphics[width=\linewidth]{dumps/sac_jepa_10k_v2/plots/latent_umap_sac_000010000_seed100001.png}
\caption{Appended SAC+JEPA endpoint deterministic inference on held-out seed
100001: terrain-potential trajectory overlay and actor-encoder UMAP.}
\label{fig:sac-jepa-trajectory-umap}
\end{figure}

\begin{figure}[t]
\centering
\includegraphics[width=\linewidth]{dumps/sac_jepa_10k_v2/plots/dynamics_history_sac_000010000_seed100001.png}
\caption{Appended SAC+JEPA physical/latent velocity, deterministic controls,
acceleration magnitudes, and signed parallel/perpendicular acceleration traces
for the same terrain-overlay inference episode.}
\label{fig:sac-jepa-dynamics}
\end{figure}

\begin{figure}[t]
\centering
\includegraphics[width=\linewidth]{dumps/sac_jepa_10k_v2/plots/pairwise_encoder_umap_5x5_seed100001.png}
\caption{Appended SAC+JEPA pairwise joint UMAP grid for checkpoints 2,048
through 10,000, evaluated on one common held-out trajectory.}
\label{fig:sac-jepa-pairwise}
\end{figure}

\begin{figure}[t]
\centering
\includegraphics[width=0.76\linewidth]{dumps/sac_jepa_10k_v2/plots/encoder_cka_matrix_seed100001.png}
\caption{Appended SAC+JEPA linear CKA matrix. Early-to-endpoint CKA is 0.53,
showing a material change in encoder geometry over this run.}
\label{fig:sac-jepa-cka}
\end{figure}

\subsection{SAC-only versus SAC+JEPA at the same 10k budget}
The prior SAC-only endpoint is $1.655\pm0.056$ while this SAC+JEPA endpoint is
$1.278\pm0.042$: a decrease of 0.377 (22.8\%) for this one matched short run.
This is not evidence against JEPA in general, but it rejects the claim that
$\lambda_{\rm JEPA}=0.1$ is already beneficial here. The next controlled
comparison should repeat seeds and test lower coefficients (0.025 and 0.05),
while retaining the complete append-only diagnostic flush for every run.

\subsection{Post-run erratum and current SAC+JEPA safeguards}
The $\lambda_{\rm JEPA}=0.1$ SAC+JEPA result above is retained as an
append-only historical result, but its training code updated the EMA target
once per replay-gradient update. With 16 environments and one update per
transition, this makes the target follow the online encoder much faster than
the nominal 0.01 update fraction suggests. It is therefore valid evidence that
the original branch ran, but not a clean test of a slowly moving predictive
target.

The current implementation corrects this. It logs the unregularized SAC actor
objective as \code{actor\_loss}, the optimized combined objective as
\code{actor\_total\_loss}, JEPA MSE as \code{jepa\_loss}, and online encoder
spread as \code{jepa\_embedding\_std}. The target now updates per 64 collected
environment transitions, independent of replay \code{updates\_per\_step}; the
new SAC-specific overlay begins at $\lambda_{\rm JEPA}=0.025$. A new,
multi-seed SAC+JEPA run is required before drawing any performance conclusion
from the corrected branch.

SAC now rejects a positive JEPA coefficient unless the SAC-specific overlay
also supplies a positive variance-floor coefficient. The variance term keeps
the online representation feature standard deviations above the configured
floor and is logged as \code{jepa\_variance\_loss}; it complements, rather
than substitutes for, the slower transition-based EMA target. This safeguard
prevents the generic PPO-oriented JEPA overlay from silently launching an
underconstrained SAC experiment.

\section{Appended run: corrected SAC+JEPA 3090 validation (September 20, 2026)}
This is a new isolated SAC+JEPA run after the metric, target-cadence, and
variance-floor corrections described above. It uses
\code{config/sac\_jepa.yaml}: $\lambda_{\rm JEPA}=0.025$, EMA update fraction
0.01 every 64 collected transitions, variance-floor coefficient 0.01, and
target standard deviation 0.1. It is therefore the appropriate SAC+JEPA
comparison, rather than the earlier 0.1 pre-safeguard run.

At 10,000 transitions the corrected branch reaches held-out mean speed
$1.403\pm0.020$. This improves the older SAC+JEPA run by 0.125, but remains
0.252 (15.2\%) below the matched SAC-only result of $1.655\pm0.056$. The JEPA
MSE decreases from $3.59\times10^{-4}$ to $1.06\times10^{-4}$ while online
encoder spread remains nonzero (0.108--0.127). The variance floor is active
only at the first checkpoint, then zero once the requested spread is reached.

\begin{figure}[t]
\centering
\includegraphics[width=\linewidth]{dumps/sac_jepa_corrected_10k/plots/learning_curve.png}
\caption{Appended corrected SAC+JEPA RTX 3090 learning curve. The JEPA loss
falls smoothly and is plotted alongside the held-out speed.}
\label{fig:sac-jepa-corrected-learning}
\end{figure}

\begin{figure}[t]
\centering
\includegraphics[width=\linewidth]{dumps/sac_jepa_corrected_10k/plots/latent_umap_sac_000010000_seed100001.png}
\caption{Appended corrected SAC+JEPA endpoint: deterministic terrain-potential
overlay and actor-encoder UMAP on held-out seed 100001.}
\label{fig:sac-jepa-corrected-umap}
\end{figure}

\begin{figure}[t]
\centering
\includegraphics[width=\linewidth]{dumps/sac_jepa_corrected_10k/plots/dynamics_history_sac_000010000_seed100001.png}
\caption{Appended corrected SAC+JEPA velocity, controls, acceleration, and
parallel/perpendicular acceleration diagnostic.}
\label{fig:sac-jepa-corrected-dynamics}
\end{figure}

\begin{figure}[t]
\centering
\includegraphics[width=\linewidth]{dumps/sac_jepa_corrected_10k/plots/pairwise_encoder_umap_5x5_seed100001.png}
\caption{Appended corrected SAC+JEPA pairwise UMAP comparison across its five
saved checkpoints.}
\label{fig:sac-jepa-corrected-pairwise}
\end{figure}

\begin{figure}[t]
\centering
\includegraphics[width=0.76\linewidth]{dumps/sac_jepa_corrected_10k/plots/encoder_cka_matrix_seed100001.png}
\caption{Appended corrected SAC+JEPA CKA matrix. Early-to-endpoint CKA is 0.61,
while consecutive late checkpoints remain more similar (0.76).}
\label{fig:sac-jepa-corrected-cka}
\end{figure}

The corrected conclusion is still negative at this budget: SAC+JEPA is now
well-conditioned and measurably non-collapsed, yet it does not beat SAC-only.
The remaining open question is whether a smaller coefficient or longer,
multi-seed training horizon changes that result; the next comparison must use
this corrected branch, not the historical 0.1 implementation.

\section{World-frame PPO and PPO+JEPA protocol (September 21, 2026)}
The previous velocity-aligned observation was found to be degenerate: its
reported local velocity is exactly $(\|\mathbf v\|,0)$, while a third state
entry separately reports the same speed. Thus the six-value state effectively
contained only speed, mass, and the prior action. More importantly, in that
frame a small lateral acceleration rotates the entire potential image while
leaving its centre fixed. This makes one-step image prediction nearly an
identity map and weakens the action-conditioned JEPA signal.

New PPO-only and PPO+JEPA runs therefore use a fixed world observation frame.
The image is aligned to world $+y$ and the state is
$(v_y,v_x,\|\mathbf v\|,m,a_{t-1,\parallel},a_{t-1,\perp})$, preserving both
velocity components. Both 5-million-transition RTX 3090 runs use the same
centered-Gaussian potential, $P_{\max}=6$ controller, 16 parallel
environments, and held-out seeds 100001--100004. The PPO+JEPA branch additionally
uses \code{config/jepa.yaml}, with one-step action-conditioned JEPA and
$\lambda_{\rm JEPA}=0.1$. Results will be appended only after final checkpoints
and the complete analysis flush are available.

For the cross-case diagnosis, the new
\code{driver.plot\_umap\_3d} tool runs deterministic inference for every
finalized checkpoint on the same held-out terrain, fits one joint three-dimensional
UMAP across all encoder states, and normalizes its three coordinates to RGB.
It writes (i) one 3-D latent UMAP per case and (ii) a 3-D physical
$(x,y,\text{speed})$ trajectory per case colored by that shared latent RGB code,
along with a row-level CSV. This makes behavioral differences and representation
differences inspectable without independently rescaling each run's UMAP.

\section{Implementation safeguards and ablations}
The current default is an \emph{unshared} actor--critic: actor and critic each
have their own CNN and 128-D MLP. The old shared encoder remains available as
the \code{agent.architecture=shared} ablation, and checkpoints saved before this
option load as shared automatically. The UMAP dump is always the actor encoder
representation. PPO now supports clipped value loss using
\code{training.value\_clip\_range=0.2}; setting it to \code{null} restores the
unclipped critic-loss ablation. Independent environment stepping is controlled
by \code{training.parallel\_env\_workers}; the default one worker preserves the
serial baseline, while larger values use a thread pool for independent NumPy
environments.

The potential acceleration is now the analytic derivative of the exact periodic
bilinear interpolant used to sample $\phi$, normalized using bilinear-cell
derivatives. Thus within each cell it is the negative gradient of the sampled
effective potential rather than a separately interpolated spectral gradient.
Tests verify this derivative against finite differences. Logged episode rows are
pre-step records: latent, state, action, potential, and acceleration all refer
to the same decision-time position. Environment \code{info} is consistently
post-step. Training asserts that generated training terrain seeds and configured
evaluation seeds are disjoint.

For checkpoint safety, normal inference loading uses PyTorch's
\code{weights\_only=True}; legacy pickle loading requires an explicit trusted
mode in code. Checkpoints are inference/experiment snapshots, not exact resume
points: environment states and in-progress rollout state are deliberately not
restored. Finally, cross-checkpoint UMAPs are qualitative only. Linear CKA is
the primary quantitative representation comparison because it is insensitive to
feature-coordinate rotations and permutations; a joint UMAP should be treated
as a shared-observation visualization, not proof of coordinate-wise agreement.

\section{\,Artifacts and reproduction}
The current tested learned network is
\code{dumps/power\_limited\_6\_10k/checkpoints/ppo\_000010000.pt}. Each \code{.pt} file contains
\code{model}, Adam \code{optimizer} state, the completed \code{steps}, resolved
\code{config}, and Torch, NumPy, and Python RNG states. It supports deterministic
evaluation but is not an exact training-resume snapshot.
Training records are in \code{dumps/power\_limited\_6\_10k/metrics/training.csv}.

The analysis additions write the following separate outputs:
\begin{itemize}[leftmargin=1.4em,itemsep=2pt]
 \item \code{dumps/power\_limited\_6\_10k/plots/learning\_curve.png}: corrected RTX 3090 speed curve;
 \item \code{dumps/power\_limited\_6\_pairwise\_10k/plots/pairwise\_encoder\_umap\_*.png}: pairwise
       joint UMAP grid across the corrected-run checkpoints;
 \item \code{dumps/power\_limited\_6\_pairwise\_10k/plots/encoder\_cka\_matrix\_*.png}: paired
       quantitative CKA comparison;
 \item \code{dumps/power\_limited\_6\_10k/plots/episode\_diagnostics\_*.png}: complete map, UMAP,
       latent-velocity, steering, and acceleration diagnostic;
 \item \code{dumps/power\_limited\_6\_10k/latents/}: raw 128-D latent CSV and NPZ dumps.
\end{itemize}
To regenerate the figures after a run, execute \code{python -m driver.plot\_training}
with the desired metrics path, then execute \code{python -m driver.analyze\_latent}
with a checkpoint and CUDA output directory.

\end{document}
